\documentclass[a4paper,11pt]{article}
\usepackage[utf8]{inputenc}
\usepackage[T1]{fontenc}
\usepackage[french]{babel}
\usepackage{lmodern}
\usepackage{amsmath,amssymb,amsthm}
\usepackage{mathrsfs}
\usepackage{enumitem}
\usepackage{multicol}

\setlist[enumerate]{label=\textbf{\textcolor{Couleur8}{\arabic*.}}, left=1em}

% Si vous voulez aussi modifier les listes enumerate imbriquées :
\setlist[enumerate,2]{label=\textcolor{Couleur8}{\alph*)}}
\setlist[enumerate,3]{label=\textcolor{Couleur8}{\roman*)}}
\setlist[enumerate,4]{label=\textcolor{Couleur8}{\Alph*)}}
\usepackage{graphicx}
\usepackage{xcolor}
\usepackage{tcolorbox}
\usepackage{geometry}
\usepackage{fancyhdr}
\usepackage{titlesec}
\usepackage{cancel}
\usepackage{ifthen}
\usepackage{tikz}
\usetikzlibrary{arrows.meta}
\usepackage{tkz-tab}
\usepackage{eurosym}

% OPTION PROF/ÉLÈVE
% Décommentez UNE des deux lignes suivantes :
\newboolean{prof}\setboolean{prof}{true}   % Version professeur (avec corrections des devoirs)
%\newboolean{prof}\setboolean{prof}{false}  % Version élève (sans corrections des devoirs)

\definecolor{Couleur1}{HTML}{4A5D7A} %Th
\definecolor{Couleur2}{HTML}{A5B5D6} %Prop
\definecolor{Couleur3}{HTML}{A8C68A} %Méthode
\definecolor{Couleur6}{HTML}{8A9CC1} %Def
\definecolor{Couleur7}{HTML}{E6B459} %Rq Voc
\definecolor{Couleur8}{HTML}{D36740} %Titres
\definecolor{correction}{HTML}{D36740} % Couleur pour les corrections
\definecolor{coopmathscorr}{HTML}{F15929} % Couleur corrections coopmaths

% Configuration des marges
\geometry{
	top=2cm,
	bottom=2cm,
	inner=1.5cm,
	outer=1.5cm,
}

% Police
\renewcommand{\familydefault}{\sfdefault}

% Compteurs
\newcounter{seancenum}
\newcounter{seancenumD}
\setcounter{seancenum}{1}
\setcounter{seancenumD}{1}

% Configuration des en-têtes et pieds de page
\pagestyle{fancy}
\fancyhf{}
\ifthenelse{\boolean{prof}}{
    \fancyhead[L]{\textcolor{Couleur8}{\textbf{Corrections Automatismes - Classe de seconde - Version Professeur}}}
}{
    \fancyhead[L]{\textcolor{Couleur8}{\textbf{Corrections Automatismes - Séance 8 - Classe de seconde}}}
}
\fancyhead[R]{\textcolor{Couleur8}{\textbf{2025-2026}}}
\fancyfoot[L]{\textcolor{Couleur8}{Mme Bertail et Mme Pinto}}
\fancyfoot[R]{\textcolor{Couleur8}{\thepage}}
\renewcommand{\headrulewidth}{1pt}
\renewcommand{\headrule}{\hbox to\headwidth{\color{Couleur8}\leaders\hrule height \headrulewidth\hfill}}

% Environnements personnalisés
\newtcolorbox{seance}[1]{
    colback=Couleur6!10,
    colframe=Couleur1!65,
    fonttitle=\bfseries\large,
    title=Correction Séance \theseancenum #1,
    left=5mm,
    right=5mm,
    top=3mm,
    bottom=3mm,
    arc=0mm,
    after=\stepcounter{seancenum}
}

\newtcolorbox{seanceD}[1]{
    colback=Couleur6!5,
    colframe=Couleur6!45,
    fonttitle=\bfseries\large,
    title=Correction Devoirs - Séance \theseancenumD #1,
    left=5mm,
    right=5mm,
    top=3mm,
    bottom=3mm,
    arc=0mm,
    after=\stepcounter{seancenumD}
}

\begin{document}

\setcounter{seancenum}{8}
\newpage
\begin{seance}{} %8

\begin{enumerate}
\item Multiplier par $0{,}75$ revient à multiplier par $1 - \dfrac{25}{100}$.
Cela revient donc à diminuer de $25\,\%$.
Ainsi, le taux d'évolution associé au coefficient multiplicateur $0{,}75$ est ${\color{correction}\boldsymbol{-25}}\,\%$.

\item Le nouveau prix est : $50 \times (1 - 0{,}2) = 50 \times 0{,}8 = {\color{correction}\boldsymbol{40}}$ €.

\item Diminuer de $25\,\%$ revient à multiplier par $0{,}75$ (coefficient multiplicateur).
Ainsi, le nouveau prix est donné par : $120 \times (1 - \dfrac{25}{100}) = 120 \times (1 - 0{,}25) = {\color{correction}\boldsymbol{120 \times (1 - \dfrac{25}{100})}}$

La bonne réponse est la réponse {\bfseries \color{correction}D}.

\item Le prix a baissé de $8$ €, ce qui correspond à $20\,\%$ de $40$ €.
Le prix du sweat a donc baissé de ${\color{correction}\boldsymbol{20}}\,\%$.

\item Le taux d'évolution $t$ est donné par la formule :
$t = \dfrac{\text{valeur finale} - \text{valeur initiale}}{\text{valeur initiale}}$

Ici : $t = \dfrac{225 - 150}{150} = \dfrac{75}{150} = 0{,}5$

Le taux d'évolution est donc de ${\color{correction}\boldsymbol{50}}\,\%$.

La bonne réponse est la réponse {\bfseries \color{correction}B}.

\item Le coefficient multiplicateur associé à une augmentation de $10\,\%$ est $1{,}1$.
Le coefficient multiplicateur associé à deux augmentations successives de $10\,\%$ est le produit des coefficients multiplicateurs, soit $1{,}1 \times 1{,}1 = 1{,}21$.
L'évolution globale est donnée par $1{,}21 - 1 = 0{,}21 = {\color{correction}\boldsymbol{+21}}\,\%$.

\end{enumerate}

\end{seance}

\end{document}